% 五、问题一
\section{问题一：确定性计划购电模型}
\subsection{建模思路}
问题一给定一天的电价、负载与光伏，要求在 0:00 一次性确定全天 144 段的购电量与储能充放动作，日初、日末储电量相等，不允许紧急购电。本文先由能量平衡、库存递推与周期约束建立线性规划，获得最优购电费；再证明其价值函数凸、分段线性且边际单调（命题~\ref{prop:convex}），并以凸折线精确表示的连续动态规划独立复算——两条路径的互证为后三问复用同一优化核与折线库奠定基础。

\subsection{确定性线性规划模型}
决策变量为各段购电量 $G_t$、充电量 $C_t$、放电量 $D_t$、富余电量 $W_t$ 与段末储电量 $S_t$，约束来自三方面物理事实，逐条推导如下。

\textbf{能量平衡}：每段供给（购电 + 光伏 + 放电）等于需求（负载 + 充电 + 富余），即
\begin{equation}
G_t+V_t+D_t=L_t+C_t+W_t,\qquad t=1,\dots,144.
\label{eq:q1balance}
\end{equation}
题面“供电不可低于负载”由 $W_t\ge0$ 表达：允许多买，多买部分作为富余弃掉。

\textbf{库存递推与界}：充电入库乘效率 $\eta$、放电出库除以 $\eta$，
\begin{equation}
S_t=S_{t-1}+\eta C_t-D_t/\eta,\qquad 1200\le S_t\le10800,
\label{eq:q1soc}
\end{equation}
且每段充放电量受功率限额约束 $0\le C_t,D_t\le\bar P=5000/6$\,kWh。

\textbf{周期库存}：题面要求 0:00 与 24:00 储电量相同，
\begin{equation}
S_0=S_{144}=6000.
\label{eq:q1periodic}
\end{equation}

目标为全天购电费最小。综合起来得到完整的线性规划模型：
\begin{equation}
\begin{aligned}
\min_{G,C,D,W,S}\quad & \sum_{t=1}^{144}p_tG_t\\
\mathrm{s.t.}\quad
& G_t+V_t+D_t=L_t+C_t+W_t, && t=1,\dots,144,\\
& S_t=S_{t-1}+\eta C_t-D_t/\eta, && t=1,\dots,144,\\
& 1200\le S_t\le 10800,\quad S_0=S_{144}=6000, && \\
& 0\le C_t,\,D_t\le\bar P,\quad G_t,\,W_t\ge 0, && t=1,\dots,144.
\end{aligned}
\label{eq:q1lp}
\end{equation}
目标中不加入 $\varepsilon(C+D+W)$ 一类经济含义之外的正则项；解的退化（同段同时充放电）在动作恢复时按等价变换消除，不影响费用。

\subsection{连续状态动态规划与价值函数性质}
线性规划给出最优值，但只给一条最优轨迹。为了回答“库存偏离该轨迹时未来费用如何变化”——这是后三问执行层的核心——定义剩余最低费用（价值函数）
\begin{equation}
F_t(e)=\min\Bigl\{\text{段 }t\text{ 起的购电费}\;\Big|\;S_{t-1}=e,\ \text{满足式\eqref{eq:q1balance}--\eqref{eq:q1periodic}}\Bigr\},
\label{eq:q1value}
\end{equation}
它满足动态规划\cite{bertsekas}的逆向递推
\begin{equation}
F_t(e)=\min_{x\in\mathcal A_t(e)}\bigl\{p_t\,[L_t-V_t+\psi(x)]^+ + F_{t+1}(e+x)\bigr\},
\qquad
F_{145}(e)=\begin{cases}0,&e=6000,\\ +\infty,&\text{其他},\end{cases}
\label{eq:q1dp}
\end{equation}
其中 $x$ 为库存增量，$\psi(x)$ 为其交流侧电量（表~\ref{tab:symbols}），可行集 $\mathcal A_t(e)$ 由功率限额与库存界给出。

\begin{proposition}[确定性价值函数的凸分段线性结构]\label{prop:convex}
按式\eqref{eq:q1dp} 定义的 $F_t(e)$，$t=1,\dots,145$，是其有效域（$[1200,10800]$ 与终端可达性约束之交）上的凸分段线性函数，且边际价值 $-\partial F_t(e)$ 关于 $e$ 单调非增；断点数至多逐段增加常数个。
\end{proposition}

\begin{proof}
对 $t$ 逆向归纳。$F_{145}$ 为单点指示函数，凸。设 $F_{t+1}$ 凸分段线性。记单段动作费用 $a_t(x)=p_t[L_t-V_t+\psi(x)]^+$。

(i) \emph{$a_t$ 凸分段线性}：$\psi(x)$ 是过原点的两段凸折线（斜率 $\eta<1/\eta$ 递增），$L_t-V_t+\psi(x)$ 仍凸；外层 $p_t[\cdot]^+$ 是非降凸函数与凸函数的复合，故 $a_t$ 凸，断点至多两个（$x=0$ 与 $[\cdot]^+$ 的零点）。

(ii) \emph{一步递推是下确界卷积}：换元 $u=e+x$，
\[
F_t(e)=\min_{u}\{\check a_t(e-u)+F_{t+1}(u)\}=(\check a_t\,\square\,F_{t+1})(e),\qquad \check a_t(z)=a_t(-z),
\]
再与功率限额、库存界对应的区间指示函数相加。凸分段线性函数的下确界卷积仍凸分段线性：其上境图是两上境图的 Minkowski 和，等价于把两条折线的斜率段按斜率归并后首尾相接，断点数为二者之和；与区间指示函数相加即定义域截断，保凸。故 $F_t$ 凸分段线性，断点增量不超过 $a_t$ 的段数（常数）。

(iii) \emph{边际单调}：凸函数的次微分关于自变量单调非减\cite{rockafellar}，故边际价值 $-\partial F_t$ 非增。归纳完成。
\end{proof}


命题~\ref{prop:convex} 的三点含义贯穿全文：(i) $F_t$ 可以用\textbf{断点--斜率序列精确表示}，逆向递推化为凸折线的下确界卷积（斜率归并排序），全程无库存网格、无插值误差；(ii) 边际价值 $-\partial F_t(e)$ 单调非增，“库存越满、追加一度电越不值钱”，这正是储能套利的边际递减结构；(iii) 给定 $F_{t+1}$，最优动作由边际价值与当段电价的比较刻画，即“放电到保留水平为止”的阈值规则——后三问的 DP 价值执行器就是该规则在不完整信息下的直接推广。

图~\ref{fig:fig_f02_value_function} 将该结构可视化：左图为四个代表时刻的 $F_t(e)$ 凸折线，右图为 12:00 的边际价值阶梯。可以看到清晨（$t=36$ 前）边际价值高——此后是早高峰放电窗口；库存越高曲线越平——富余库存的机会成本递减。

\begin{figure}[!htbp]\centering
\safeincludegraphics[width=0.92\textwidth]{fig_f02_value_function}
\caption{问题一价值函数 $F_t(e)$（左）与 12:00 边际价值 $-\partial F_t$（右）：凸、分段线性、边际单调，验证命题~\ref{prop:convex}}
\label{fig:fig_f02_value_function}
\end{figure}

\subsection{双方法互证与求解验证}
先说明二者为何必须相等。式\eqref{eq:q1value} 按定义就是"从状态 $(t,e)$ 出发的原问题最优值"，故 $F_1(6000)$ 与线性规划\eqref{eq:q1lp} 是\emph{同一个}优化问题：可行域相同（能量平衡、库存递推与界、周期约束逐条对应），目标相同。而逆向递推式\eqref{eq:q1dp} 与定义式\eqref{eq:q1value} 的等价即 Bellman 最优性原理\cite{bertsekas}——本问题逐段费用可加、状态转移只依赖当前状态与动作，原理的适用条件全部满足；且凸折线表示对 $F_t$ 是精确的（命题~\ref{prop:convex}），递推中不引入任何离散化或截断近似。因此\textbf{理论上 LP 最优值恒等于 $F_1(6000)$，任何数值差异只可能来自实现错误}——这正是对照的意义：它检验的不是模型，而是两条独立实现路径的代码正确性。

线性规划用 SciPy\cite{scipy} 的 HiGHS 求解器求解；连续动态规划按式\eqref{eq:q1dp} 以凸折线精确递推并正向恢复动作。两条完全独立的实现路径给出同一最优值：
\[
\text{LP}=35126.9486\ \text{元},\qquad F_0(6000)=35126.9486\ \text{元},\qquad \text{差}=7.3\times10^{-12}\ \text{元},
\]
再对非网格初值、零电价、单点终端等 750 组变式做同样对照，最大费用差 $3.64\times10^{-12}$ 元。互证排除了实现错误，也使后三问放心复用同一优化核与折线库。

\subsection{最优策略及经济解读}
最优计划全天购电 59482.70\,kWh、购电费 \textbf{35126.95 元}。作为对照，若不使用储能、逐段按需购电，全天费用为 48052.05 元：储能的峰谷套利把日购电费降低 12925.10 元（\textbf{26.90\%}）。指定时段结果见表~\ref{tab:q1purchase} 与表~\ref{tab:q1storage}，完整策略已存入 result1.xlsx。

\begin{table}[!htbp]
\centering\small
\caption{问题一指定时段购电量（kWh）}\label{tab:q1purchase}
\begin{tabular}{lrrrrrr}
\toprule
时段起点 & 10:00 & 12:00 & 14:00 & 16:00 & 18:00 & 20:00 \\
\midrule
购电量 & 0.00 & 480.41 & 0.00 & 445.43 & 531.89 & 0.00 \\
\bottomrule
\end{tabular}
\end{table}

\begin{table}[!htbp]
\centering\small
\caption{问题一四小时充放电汇总（kWh）}\label{tab:q1storage}
\begin{tabular}{lrr@{\qquad}lrr}
\toprule
时段 & 充电量 & 放电量 & 时段 & 充电量 & 放电量 \\
\midrule
0:00--4:00  & 4500.00 &    0.00 & 4:00--8:00   &  833.33 & 6365.84 \\
8:00--12:00 & 4787.96 & 1703.00 & 12:00--16:00 & 5286.04 &   91.10 \\
16:00--20:00 &   0.00 & 5780.13 & 20:00--24:00 & 5333.33 & 2859.87 \\
\bottomrule
\end{tabular}
\end{table}


\begin{figure}[!htbp]\centering
\safeincludegraphics[width=0.9\textwidth]{fig_f03_q1_day}
\caption{问题一最优调度：(a) 负载、光伏与计划购电功率；(b) 储电量轨迹。低价谷段（0--6 时、12--14 时）满功率充电、购电顶到负载之上，高价峰段（7--11 时、17--20 时）放电、购电归零}
\label{fig:fig_f03_q1_day}
\end{figure}

图~\ref{fig:fig_f03_q1_day} 显示最优策略的经济机理：购电集中在电价谷段并同步充电，峰段完全依靠光伏与放电覆盖负载；储电量在两个峰前充至接近上限、峰内放至接近下限，一天完成约两个完整充放循环——在效率损耗 $\eta^2=0.81$ 下，只有峰谷价差超过约 $1/0.81-1\approx23\%$ 的时段对才值得套利，这解释了为何部分中等价差时段储能保持静默。周期约束的对偶价格给出 24:00 库存的边际影子价，与次日 0:00 的 $F$ 曲线斜率一致，说明周期解自洽。
